% =====================================================================
% BACKUP: Full encoder (BERT) CVT descriptor details.
% Moved out of thesis_report/sections/09b_detailed_pipeline.tex when the
% CVT subsection was compacted. Use this to build the "Encoder Pruning
% Descriptors" section of the Search Descriptors appendix
% (sec:appendix_descriptors). Self-contained content snippet: paste the
% body into the appendix source and adjust \section/\subsection levels.
%
% Notation: a_l(x) = active attention heads in layer l;
%           n_l(x) = active feed-forward neurons in layer l;
%           headcost, neuroncost = per-head / per-neuron operation costs;
%           nlayers = L (number of layers); nffn = dense FFN width.
% =====================================================================

\section*{Encoder Pruning Descriptors}

Each candidate mask $x$ is mapped to a six-coordinate descriptor
\[
z(x)=(z_1(x),\ldots,z_6(x))\in[0,1]^6 ,
\]
whose coordinates capture the allocation properties that determine whether
two masks should be refined by the same local mutation process.

\begin{table}[H]
\centering
\footnotesize
\setlength{\tabcolsep}{4pt}
\begin{tabularx}{0.92\textwidth}{c l X}
\toprule
Coordinate & Quantity & Interpretation \\
\midrule
$z_1$ & attention MAC share & separates head-heavy from neuron-heavy allocations \\
$z_2$ & surviving feed-forward width & measures how much dense feed-forward capacity remains \\
$z_3$ & complete active-layer fraction & measures how much depth retains both branches \\
$z_4$ & budget entropy across layers & distinguishes spread-out and concentrated allocations \\
$z_5$ & head centre of mass & locates attention capacity along depth \\
$z_6$ & neuron centre of mass & locates feed-forward capacity along depth \\
\bottomrule
\end{tabularx}
\caption{Six descriptor coordinates used to partition the feasible sweep pool.}
\label{tab:cvt_descriptor}
\end{table}

The descriptor begins with the two global counts induced by the mask. For
candidate $x$, $\actl(x)$ is the number of active attention heads in layer
$\ell$, and $\nactl(x)$ is the number of active feed-forward neurons in that
layer. Summing these counts over all layers gives the total surviving
attention and feed-forward capacity:
\[
H_{\mathrm{tot}}(x)=\sum_{\ell}\actl(x),
\qquad
N_{\mathrm{tot}}(x)=\sum_{\ell}\nactl(x).
\]
These totals define the first view of the candidate: its balance between
attention and feed-forward capacity. The constants $\headcost$ and
$\neuroncost$ are the operation costs of one active head and one active
feed-forward neuron from the budget model. The coordinate $z_1$ is the
fraction of the mask-dependent cost assigned to attention, and $z_2$ is the
surviving feed-forward width normalised by the dense feed-forward width
$\nlayers\nffn$:
\[
z_1(x)
=
\frac{H_{\mathrm{tot}}(x)\headcost}
{H_{\mathrm{tot}}(x)\headcost+N_{\mathrm{tot}}(x)\neuroncost},
\qquad
z_2(x)
=
\frac{N_{\mathrm{tot}}(x)}{\nlayers\nffn}.
\]
The next view describes how the surviving computation is distributed across
depth. The coordinate $z_3$ measures how many layers remain complete, meaning
that both the attention branch and the feed-forward branch are still active:
\[
z_3(x)
=
\frac{\left|\{\ell:\actl(x)>0 \wedge \nactl(x)>0\}\right|}{\nlayers}.
\]
The coordinate $z_4$ measures whether the candidate spreads its budget over
many layers or concentrates it in a small number of layers. Here $b_\ell$ is
the operation cost assigned to layer $\ell$, and $p_\ell$ is the fraction of
the candidate's layer-wise budget carried by that layer:
\[
z_4(x)
=
-\frac{1}{\log \nlayers}
\sum_{\ell:b_\ell>0}p_\ell\log p_\ell,
\qquad
p_\ell=\frac{b_\ell}{\sum_k b_k},
\qquad
b_\ell=\actl(x)\headcost+\nactl(x)\neuroncost .
\]
The entropy coordinate $z_4$ is large when the budget is distributed across
depth and small when it is concentrated. The final view records where the
surviving capacity is located. The layer index $\ell$ is a depth coordinate;
division by $\nlayers-1$ normalises the weighted average into the unit
interval. The coordinate $z_5$ is the centre of mass of active heads, and
$z_6$ is the centre of mass of active feed-forward neurons:
\[
z_5(x)
=
\frac{1}{(\nlayers-1)H_{\mathrm{tot}}(x)}
\sum_{\ell}\ell\,\actl(x),
\qquad
z_6(x)
=
\frac{1}{(\nlayers-1)N_{\mathrm{tot}}(x)}
\sum_{\ell}\ell\,\nactl(x).
\]
Low values of $z_5$ or $z_6$ indicate capacity placed near the lower layers,
and high values indicate capacity placed near the upper layers. Taken
together, the six coordinates place masks with similar allocation balance and
similar depth distribution close to one another.

\subsection*{Standardisation}

Before distance-based clustering, the descriptor coordinates are standardised
because their empirical spreads can differ across the sweep pool:
\[
\tilde z_i=\frac{z(x_i)-\bar z}{\sigma_z},
\]
where $z(x_i)$ is the descriptor of sweep candidate $x_i$, $\bar z$ is the
componentwise sample mean, $\sigma_z$ is the componentwise sample standard
deviation, and $\tilde z_i$ is the standardised descriptor used for
clustering. Standardisation places all six coordinates on a common numerical
scale. The standardised descriptors are then partitioned by the Centroidal
Voronoi Tessellation (Lloyd fitting given in the CVT appendix,
\texttt{sec:appendix\_cvt}).
